\chapter{Introduction}

\section{Overview}
The rapid digitization of healthcare information systems has fundamentally transformed the manner in which patient diagnostics, treatment records, and medical histories are managed and shared among medical professionals and institutions. Electronic Health Records (EHRs) offer tremendous advantages in terms of accessibility, efficiency, and continuity of care. However, the centralization of highly sensitive patient data introduces profound security and privacy challenges. Traditional monolithic database architectures present a single point of failure; an unauthorized breach or a successful cyberattack can compromise millions of confidential medical records simultaneously. Consequently, safeguarding these data assets while ensuring their availability to authorized personnel is a critical concern in modern healthcare informatics. 

Blockchain technology has emerged as a promising paradigm to address these vulnerabilities by providing a decentralized, immutable, and transparent ledger for access logs and auditing. Yet, storing raw medical data directly on a public or even a permissioned blockchain is impractical due to block size limitations, high latency, and stringent privacy regulations such as HIPAA and GDPR \cite{zhang2021security, shorfuzzaman2021blockchain}. Therefore, a hybrid approach combining off-chain storage with on-chain metadata management is essential for a scalable healthcare system \cite{sharma2025hybrid, rahman2019secure}. 

This project goes a step further by recognizing that medical records are not monolithic entities. A patient may wish to grant a specialist access to their cardiology reports while restricting access to their psychiatric evaluations. This paradigm is referred to as "condition-wise" access control. achieving fine-grained, condition-specific security requires cryptographic techniques capable of compartmentalizing information efficiently.

\section{Motivation}
The primary motivation behind this endeavor is the lack of existing frameworks that integrate blockchain-based auditability with robust, fine-grained, condition-wise data encryption. Current solutions frequently utilize a single symmetric key per patient record or rely heavily on the continuous availability of an external Trusted Authority (TA). This dependency leads to two massive issues. Firstly, if the TA goes offline, the records become inaccessible. Secondly, if a single encryption key is leaked, the entirety of the patient's records is exposed. 

By employing advanced secret sharing algorithms—specifically Shamir's Secret Sharing—we can distribute cryptographic key material across a decentralized network of peer nodes \cite{li2023privacy}. This eliminates the vulnerability of a single point of failure and drastically mitigates the risks associated with Key Escrow. Our motivation is to engineer a healthcare data management ecosystem that dynamically assesses the criticality of an uploaded medical document, determines a fault-tolerance threshold, and encrypts the document strictly corresponding to its specific medical condition \cite{zhang2021priority}.

\section{Problem Statement}
The continuous rise in data breaches targeting healthcare institutions demonstrates that conventional perimeter-based security measures are inadequate \cite{jin2019review}. When patient data is stored centrally and encrypted with coarse-grained methodologies, single-point authorization mechanisms can be bypassed. The specific problem addressed in this project is: \textit{How can we design and implement a distributed, tamper-proof healthcare information system that guarantees condition-wise encryption, dynamically adapts its threshold cryptography according to data criticality using Large Language Models (LLMs) \cite{wang2024large, sakai2025large}, and leverages a permissioned blockchain for verifiable tracking without creating a single point of compromise \cite{kumar2022permissioned}?}

\section{Research Objectives}
To address the aforementioned problem, this project outlines the following core objectives:
\begin{enumerate}
    \item \textbf{Decentralized Architecture Formulation:} To design a system where the Trusted Authority performs cryptographic operations outside the blockchain ordering service boundary, preserving defense in depth.
    \item \textbf{Condition-Wise Encryption:} To implement an AES-GCM based encryption schema that generates dedicated sub-keys for specific medical conditions (e.g., Asthma, Hypertension), preventing cross-condition authorization leakage.
    \item \textbf{Threshold Cryptography Integration:} To deploy Shamir's Secret Sharing to split encryption keys into shares, distributing them securely to Fabric peer nodes utilizing Node Master Keys (NMK).
    \item \textbf{Dynamic Priority Classification:} To integrate an LLM interface that evaluates the content of a medical report and assigns a priority level (LOW, MEDIUM, HIGH), which inversely dictates the number of shares (threshold $k$) required to reconstruct the data.
    \item \textbf{Verifiable Auditing:} To maintain an immutable tracking log of all data uploads, updates, and accesses using the Hyperledger Fabric blockchain, ensuring strict accountability.
\end{enumerate}

\section{Organization of the Report}
The remainder of this report is organized as follows: Chapter 2 discusses the foundational background and literature review. Chapter 3 elaborates on the proposed system architecture, detailing the integration of Hyperledger Fabric and the Trusted Authority. Chapter 4 rigorously formulates the cryptographic methodologies and algorithms used. Chapter 5 covers the concrete implementation specifics involving Python, Node.js and Streamlit. Chapter 6 provides an exhaustive evaluation of system performance, latency breakdown, and fault tolerance. Finally, Chapter 7 concludes the report and describes potential avenues for future enhancements.
